% GENERATED FILE. Edit canonical lessons, then run npm run generate:books.
% canonical-source-hash: fnv1a64:a050da75d525dd14
% canonical-lessons: UR-C09-bhai, UR-C09-bahan, UR-C09-dost, UR-C09-khandan

\chapter{People You'd Introduce}
\label{ch:ur-people}
\hlchaptermodality{9}

\begin{chapteropening}
\textbf{By the end of this chapter:} \emph{I can name my brother, sister, friend and family in Urdu, and choose merā or merī correctly for each noun's grammatical gender.}

Say merā bhāī, merī bahan, merā dost and merā khāndān with merā/merī chosen correctly for each noun, and sort the four into inherited bhāī/bahan against Persian dost/khāndān — the same native-versus-borrowed split Chapter 8's verbs drew.
\end{chapteropening}

\section[bhāī]{\ur{بھائی} — \textquotedblleft{}brother,\textquotedblright{} and the first real English cousin among these nouns}

\label{lesson:UR-C09-bhai}



\begin{quote}
Say \textbf{mujhe paṛhnā pasand hai}. That closed the verb chapters. This chapter opens a new kind of word: nouns for the people you would actually introduce.
\end{quote}

\subsection*{You'll want to know first — one word}
\begin{quote}
\textbf{\ur{بھائی}} — \emph{bhāī} — \textbf{brother}
\end{quote}

\begin{quote}
\textbf{\ur{میرا} \ur{بھائی۔}} — \emph{merā bhāī.} — \textquotedblleft{}{[}He's{]} my brother.\textquotedblright{}
\end{quote}

Urdu, like English, can drop the subject when it is obvious from context — pointing at someone while you say this is enough.

\subsection*{The letters in this word}
From the right edge: \textbf{\ur{ب}} \emph{b}, then \textbf{\ur{ھ}} aspirating it to \emph{bh}, then \textbf{\ur{ا}} long \emph{ā}, then \textbf{\ur{ئ}}, then \textbf{\ur{ی}} \emph{ī}.

\textbf{\ur{ئ}} is new: \textbf{ye with hamza above}. It is a fourth shape in the \emph{ye} family, alongside consonantal \emph{y} in \textbf{\ur{کیا}}, long \emph{ī} in \textbf{\ur{جی}}, and broad final \emph{e} in \textbf{\ur{کیسے}}. Its job is narrower than any of those: it marks a glide between two vowels that would otherwise collide — here, between \textbf{ā} and \textbf{ī}. \textbf{\ur{بھ}} is new too, extending the two-eyed \emph{he} you already know from \textbf{\ur{جھ}}, \textbf{\ur{ٹھ}}, \textbf{\ur{ڑھ}} and \textbf{\ur{چھ}} to a fifth base letter, \textbf{\ur{ب}}.

\begin{grammarlens}[title={Grammar Lens: a noun with a gender}]
\textbf{\ur{بھائی}} is grammatically \textbf{masculine} — the same class as \textbf{\ur{نام}} in \textbf{merā nām}. That is why \textbf{merā} needs no change here. Not every noun agrees this easily; the next lesson breaks the pattern.
\end{grammarlens}

\begin{cousinweb}
\textbf{bhāī} comes from Sanskrit \emph{bhr\'{\={a}}tṛ}, \textquotedblleft{}brother,\textquotedblright{} from Proto-Indo-Iranian \textbf{*b\textsuperscript{h}ráHtā}, from Proto-Indo-European \textbf{*b\textsuperscript{h}réh\textsubscript{2}tēr}. English \textbf{brother} comes from the very same root. Unlike \textbf{\ur{پسند}}'s disputed link to \emph{candēre} two lessons ago, this one is settled: Sanskrit and English are both Indo-European, so \emph{bhāī} and \emph{brother} are real relatives, not two languages landing on the same idea by accident.
\end{cousinweb}

\subsection*{Your turn}
\begin{itemize}
\raggedright
  \item \emph{Say it:} \textbf{bhāī} — brother; then \textbf{merā bhāī}
  \item \emph{Name:} the fourth \emph{ye} shape — hamza above, between two vowels
  \item \emph{Connect:} \textbf{bhāī} $\leftarrow$ \emph{bhr\'{\={a}}tṛ} $\to$ English \textbf{brother}, real cousins
  \item \emph{Say it:} \textbf{mujhe paṛhnā pasand hai} once more, then \textbf{merā bhāī}
\end{itemize}

\begin{tcolorbox}[breakable,colback=teal!4,colframe=teal!35!black,title={Before you move on}]
What does \textbf{\ur{بھائی}} mean, and is \textbf{\ur{بھائی}} masculine or feminine? (\textbf{Brother}; \textbf{masculine}.) Which English word is its real cousin, not just a resemblance? (\textbf{Brother}.) Name the four jobs of the \emph{ye} family now. (\textbf{Consonantal \emph{y}, long \emph{ī}, final \emph{e}, and hamza-glide.})

Sources: \href{https://en.wiktionary.org/wiki/\%D8\%A8\%DA\%BE\%D8\%A7\%D8\%A6\%DB\%8C}{Wiktionary: \ur{بھائی}}.
\end{tcolorbox}
\section[bahan]{\ur{بہن} — \textquotedblleft{}sister,\textquotedblright{} and the word that makes \ur{میرا} move}

\label{lesson:UR-C09-bahan}



\begin{quote}
Say \textbf{merā bhāī} — my brother, a real cousin of the English word. Here is the other half of the pair, and it changes something you have used unchanged since Chapter 2.
\end{quote}

\subsection*{You'll want to know first — one word}
\begin{quote}
\textbf{\ur{بہن}} — \emph{bahan} — \textbf{sister}
\end{quote}

\begin{quote}
\textbf{\ur{میری} \ur{بہن۔}} — \emph{merī bahan.} — \textquotedblleft{}{[}She's{]} my sister.\textquotedblright{}
\end{quote}

Not \textbf{\ur{میرا} \ur{بہن}}. Read on for why.

\begin{grammarlens}[title={Grammar Lens: \ur{میرا} bends to the noun, not to you}]
\textbf{\ur{بہن}} is grammatically \textbf{feminine}. Urdu's word for \textquotedblleft{}my\textquotedblright{} carries that agreement: \textbf{\ur{میرا}} \emph{merā} before a masculine noun, \textbf{\ur{میری}} \emph{merī} before a feminine one — regardless of whether the speaker is a man or a woman. It is the noun's gender that decides, never the speaker's.

\begin{quote}
\textbf{\ur{میرا} \ur{بھائی،} \ur{میری} \ur{بہن۔}} — \emph{merā bhāī, merī bahan.} — my brother, my sister.
\end{quote}

\textbf{\ur{نام}} in Chapter 2's \textbf{merā nām} is masculine, which is why that phrase never had to teach you this rule — one gender, hidden in plain sight, until a feminine noun finally showed up.
\end{grammarlens}

\begin{cousinweb}
\textbf{bahan} continues Sanskrit \emph{bhaginī}, traditionally connected to \textbf{bhaj-}, \textquotedblleft{}to share, to portion out\textquotedblright{} — a sister pictured as a co-heir, sharing the family's portion. Specialists flag a real problem with the sound history, though: the aspiration in \emph{bhaginī} does not shift the way this derivation needs it to. Hold this one loosely, the way \textbf{\ur{پسند}}'s \emph{candēre} link was held loosely — a plausible story, not a settled one like \emph{bhāī}'s.
\end{cousinweb}

\subsection*{Your turn}
\begin{itemize}
\raggedright
  \item \emph{Say it:} \textbf{bahan} — sister; then \textbf{merī bahan}
  \item \emph{Contrast:} \textbf{merā bhāī} — \textbf{merī bahan}
  \item \emph{Say it:} why \textbf{merā} changed — the noun's gender, not the speaker's
  \item \emph{Say it:} \textbf{bahan} $\leftarrow$ \emph{bhaginī} $\leftarrow$ \emph{bhaj-}, \textquotedblleft{}to share\textquotedblright{} — held loosely
\end{itemize}

\begin{tcolorbox}[breakable,colback=teal!4,colframe=teal!35!black,title={Before you move on}]
Why is it \textbf{\ur{میری} \ur{بہن}} and not \textbf{\ur{میرا} \ur{بہن}}? (\textbf{\ur{بہن} is feminine, and merā/merī agrees with the noun's gender.}) What does \emph{bhaj-} mean, and is the connection to \emph{bahan} certain? (\textbf{\textquotedblleft{}To share\textquotedblright{}; no, it is debated.})

Sources: \href{https://en.wiktionary.org/wiki/\%D8\%A8\%DB\%81\%D9\%86}{Wiktionary: \ur{بہن}}.
\end{tcolorbox}
\section[dost]{\ur{دوست} — \textquotedblleft{}friend,\textquotedblright{} and the relationship \ur{تم} was made for}

\label{lesson:UR-C09-dost}



\begin{quote}
Back in Chapter 3 you learned three Urdu words for \textquotedblleft{}you\textquotedblright{} — respectful \textbf{āp}, familiar \textbf{tum}, intimate \textbf{tū} — and were told to use \textbf{tum} once a relationship supports familiarity. Here is the relationship that phrase meant.
\end{quote}

\subsection*{You'll want to know first — one word}
\begin{quote}
\textbf{\ur{دوست}} — \emph{dost} — \textbf{friend}
\end{quote}

\begin{quote}
\textbf{\ur{میرا} \ur{دوست۔}} — \emph{merā dost.} — \textquotedblleft{}{[}He's/She's{]} my friend.\textquotedblright{}
\end{quote}

\begin{grammarlens}[title={Grammar Lens: a noun that will not pick a gender}]
\textbf{\ur{بھائی}} was masculine and needed \textbf{merā}. \textbf{\ur{بہن}} was feminine and needed \textbf{merī}. \textbf{\ur{دوست}} breaks the pattern on purpose: it takes \textbf{the same form} for a male or a female friend, and \textbf{merā dost} works for either. Native Urdu nouns almost always commit to a gender; this Persian loan simply does not carry the machinery that would make it agree.
\end{grammarlens}

\begin{cousinweb}
\textbf{dost} is Persian, inherited from Old Persian \textbf{dauštā}, from Proto-Iranian \textbf{*jawšt\'{\={a}}}, from Proto-Indo-European \textbf{*ǵews-}, \textquotedblleft{}to taste, to choose, to enjoy.\textquotedblright{} That root also gave English \textbf{choose} and Latin \emph{gustus} — behind English \textbf{gusto}. A friend, on this history, is someone \textbf{tasted and chosen}, the same picture as \emph{lenā}'s worn-down \textquotedblleft{}take.\textquotedblright{}

Persian's own word for \textquotedblleft{}ask,\textquotedblright{} \textbf{porsīdan}, turned out to share one root with Urdu's inherited \textbf{\ur{پوچھنا}}. \textbf{\ur{دوست}} is not that kind of cousin — Urdu has no inherited word from \emph{*ǵews-} to compare it against — but it is the same general shape: an old Indo-European root, arriving in Urdu by the Persian road alone.
\end{cousinweb}

\begin{grammarlens}[title={Grammar Lens: where \ur{تم} (\emph{tum}) actually lives}]
Chapter 3 taught \textbf{āp / tum / tū} and told you to move to \textbf{tum} once familiarity is safe. A \textbf{\ur{دوست}} is exactly that relationship: someone you have moved off \textbf{āp} for. You do not need a new rule — just the word for the person the old rule was describing.
\end{grammarlens}

\subsection*{Your turn}
\begin{itemize}
\raggedright
  \item \emph{Say it:} \textbf{dost} — friend; then \textbf{merā dost}, for either a man or a woman
  \item \emph{Contrast:} \textbf{merā bhāī}, \textbf{merī bahan}, \textbf{merā dost} — two that agree, one that does not
  \item \emph{Connect:} \textbf{dost} $\leftarrow$ \emph{*ǵews-} $\to$ English \textbf{choose}, \textbf{gusto}
  \item \emph{Say it:} \textbf{āp, tum, tū} — then name which one fits a \textbf{\ur{دوست}}
\end{itemize}

\begin{tcolorbox}[breakable,colback=teal!4,colframe=teal!35!black,title={Before you move on}]
Does \textbf{\ur{دوست}} change form for a male or female friend? (\textbf{No — it stays \ur{دوست} either way.}) Which two English words share \textbf{\ur{دوست}}'s root? (\textbf{Choose} and \textbf{gusto}.) Which of \textbf{āp / tum / tū} fits a \textbf{\ur{دوست}}? (\textbf{tum}.)

Sources: \href{https://en.wiktionary.org/wiki/\%D8\%AF\%D9\%88\%D8\%B3\%D8\%AA}{Wiktionary: \ur{دوست}}.
\end{tcolorbox}
\section[khāndān]{\ur{خاندان} — \textquotedblleft{}family,\textquotedblright{} the word your three new nouns never needed}

\label{lesson:UR-C09-khandan}



\begin{quote}
Three people-words so far: \textbf{bhāī}, \textbf{bahan}, \textbf{dost} — a specific brother, a specific sister, a specific friend. Not one of them names the group they belong to.
\end{quote}

\subsection*{You'll want to know first — one word}
\begin{quote}
\textbf{\ur{خاندان}} — \emph{khāndān} — \textbf{family}
\end{quote}

\begin{quote}
\textbf{\ur{میرا} \ur{خاندان۔}} — \emph{merā khāndān.} — \textquotedblleft{}{[}This is{]} my family.\textquotedblright{}
\end{quote}

\textbf{\ur{خاندان}} is grammatically masculine, so \textbf{merā} stays unchanged — the same class as \textbf{\ur{بھائی}}.

\begin{cousinweb}
\textbf{\ur{خاندان}} is Classical Persian, borrowed into Urdu whole. Unlike \textbf{\ur{خدا} \ur{حافظ}}, whose Persian \emph{khudā} and Arabic-rooted \emph{hāfiz} this book could pull apart piece by piece, no source here commits to what \textbf{\ur{خاندان}} is built from internally — so this one stays a single unanalyzed loan, not a compound to take apart.

Line up this chapter's four nouns by where they came from. \textbf{\ur{بھائی}} and \textbf{\ur{بہن}} are inherited, straight from Sanskrit, the same road as every kin word before them. \textbf{\ur{دوست}} and \textbf{\ur{خاندان}} are Persian, the same road as \textbf{\ur{خدا}} and \textbf{\ur{حافظ}} back in Chapter 5. The people you would point to and name stayed native; the words for choosing them as friends and grouping them as family reached outside the language.
\end{cousinweb}

\subsection*{Your turn}
\begin{itemize}
\raggedright
  \item \emph{Say it:} \textbf{khāndān} — family; then \textbf{merā khāndān}
  \item \emph{Say it:} all four — \textbf{bhāī, bahan, dost, khāndān}
  \item \emph{Sort:} inherited — \textbf{bhāī, bahan}; Persian — \textbf{dost, khāndān}
  \item \emph{Choose:} \textbf{merā} or \textbf{merī} for each of the four
\end{itemize}

\begin{tcolorbox}[breakable,colback=teal!4,colframe=teal!35!black,title={Before you move on}]
Which two of this chapter's nouns are inherited from Sanskrit, and which two are Persian loans? (\textbf{\ur{بھائی}, \ur{بہن} inherited; \ur{دوست}, \ur{خاندان} Persian.}) Say all four people-words with \textbf{merā} or \textbf{merī} correctly chosen for each.

Sources: \href{https://en.wiktionary.org/wiki/\%D8\%AE\%D8\%A7\%D9\%86\%D8\%AF\%D8\%A7\%D9\%86}{Wiktionary: \ur{خاندان}}.
\end{tcolorbox}
